\documentclass[11pt]{article}
\usepackage{fltcourse}

\title{Level 9: Deformation theory and the automorphy lifting theorem}
\author{Kevin Buzzard}
\date{}

\begin{document}
\maketitle

\noindent
We have reduced Fermat's Last Theorem to the three statements $B_{8a}$
(automorphic lifting), $B_{8b}$ (automorphic spreading out) and $B_{8c}$
(potential automorphy). Each of these is out of reach by soft means alone: as
we noted, the automorphic analogues of lifting and spreading out are easy, but
they only control a representation after restriction to a totally real field.
The instrument that converts control of $\rho|_{\GF}$ into control of $\rho$ is
an \emph{automorphy lifting theorem}, which identifies a Galois deformation
ring with a Hecke algebra. This level develops the deformation theory the
statement rests on, states the automorphy lifting theorem, and folds it into a
new boss statement
\[
  B_9 := B_{8a}\wedge B_{8b}\wedge B_{8c}\wedge \text{(ALT)}.
\]
The boss theorem itself is then immediate; the real work --- using the
automorphy lifting theorem to kill $B_{8a}$, $B_{8b}$, $B_{8c}$ one at a time
--- is the business of the final three levels.

\section{Deformation functors}

The deformation theory below is due to Mazur (1989). We recall it briefly; the
category-theoretic language is standard, and the constructions are formalized.

Fix a finite field $k$ of characteristic $\ell$. Let $\calC_k$ be the category
whose objects are coefficient rings $R$ equipped with an isomorphism
$R/\m_R\xrightarrow{\sim}k$, and whose morphisms are continuous local
homomorphisms compatible with the maps to $k$. (Working with a fixed
identification of the residue field with $k$, rather than with ``rings with
residue field $k$'', avoids a set-theoretic nuisance.) The initial object is
$\Zl$ or, more generally, the Witt vectors $W(k)$; the residue field $k$ is
terminal.

Let $G$ be a profinite group, $V$ a finite-dimensional $k$-vector space, and
$\rhobar\colon G\to\Aut_k(V)$ a continuous representation. The \emph{deformation
functor} $D_{\rhobar}\colon\calC_k\to\mathrm{Set}$ sends $(R,i)$ to the set of
isomorphism classes of pairs $(W,\rho)$, where $W$ is a finite free $R$-module,
$\rho\colon G\to\Aut_R(W)$ is continuous, and there is a $G$-equivariant
isomorphism $W\otimes_{R,i}k\cong V$; a morphism $\phi\colon R\to R'$ acts by
base change $(W,\rho)\mapsto(W\otimes_{R,\phi}R',\rho\otimes R')$. A functor
$\calC_k\to\mathrm{Set}$ isomorphic to $R\mapsto\Hom_{\calC_k}(R^{\univ},R)$ is
\emph{representable}, by $R^{\univ}$.

\begin{theorem}[Mazur, 1989]\label{thm:mazur}
Let $K$ be a number field, $S$ a finite set of finite places, $K^S$ the maximal
extension of $K$ unramified outside $S$, and $G=\Gal(K^S/K)$. If
$\rhobar\colon G\to\Aut_k(V)$ is absolutely irreducible then $D_{\rhobar}$ is
representable, by a \emph{universal deformation ring} $R^{\univ}_{\rhobar}$, with
a universal deformation
$\rho^{\univ}\colon G\to\GL(W^{\univ})$ through which every deformation
uniquely factors.
\end{theorem}

The hypothesis on $G$ is the Hermite--Minkowski finiteness: $K$ has only finitely
many extensions of bounded degree unramified outside $S$. Mazur's proof uses
Schlessinger's criterion; there are also elementary proofs (Faltings, de
Smit--Lenstra). We use the theorem as stated.

\section{Deformation conditions}

We shall want to deform not all the way, but subject to conditions --- ``flat
at $\ell$'', ``cyclotomic determinant'', ``hardly ramified''. A
\emph{deformation condition} is an isomorphism-invariant predicate $P$ on pairs
$(W,\rho)\in D_{\rhobar}(R)$, defined for all $R\in\calC_k$; it cuts out a
subset $P_{\rhobar}(R)\subseteq D_{\rhobar}(R)$. The following criterion, due to
de Smit--Lenstra and Ramakrishna, tells us when it cuts out a
\emph{representable} subfunctor, and this is the one place in the level where we
prove something.

\begin{theorem}[de Smit--Lenstra, Ramakrishna]\label{thm:condition}
With $G$, $k$, $\rhobar$ as in Theorem~\ref{thm:mazur}, let $P$ be a deformation
condition satisfying:
\begin{enumerate}
\item[\textup{(1)}] $(W,\rho)$ satisfies $P$ iff $W/JW$ satisfies $P$ for every
open ideal $J\subseteq R$;
\item[\textup{(2)}] if $I,J\subseteq R$ are open and $W/IW$, $W/JW$ satisfy
$P$, then $W/(I\cap J)W$ satisfies $P$;
\item[\textup{(3)}] for an inclusion $R\subseteq R'$ of finite objects of
$\calC_k$ and $W$ over $R$, $W$ satisfies $P$ iff $W\otimes_R R'$ does.
\end{enumerate}
Then $P_{\rhobar}$ is representable by $R^{\univ}_{\rhobar}/J$ for a closed ideal
$J$.
\end{theorem}

\begin{proof}
For a deformation $W$ over $R$, let $J_W$ be the intersection of the open ideals
$I$ with $W/IW$ satisfying $P$. By (2) this family is closed under finite
intersection, so $R/J_W=\varprojlim_I R/I$ over it, and by (1) $W/J_WW$ satisfies
$P$; moreover $W/IW$ satisfies $P$ iff $J_W\subseteq I$, and (again by (1), as
$R$ is profinite) $W$ itself satisfies $P$ iff $J_W=0$.

We claim $J_W$ is compatible with base change: for a morphism $\phi\colon R\to R'$
and $W'=W\otimes_{R,\phi}R'$, we have $\phi(J_W)R'=J_{W'}$. Indeed, for an open
ideal $I'\subseteq R'$ with preimage $I=\phi^{-1}(I')$, the map
$R/I\hookrightarrow R'/I'$ is an injection of finite objects, so by (3),
$W/IW$ satisfies $P$ iff $W'/I'W'$ does; that is, $J_W\subseteq I=\phi^{-1}(I')$
iff $J_{W'}\subseteq I'$, i.e.\ $\phi(J_W)R'\subseteq I'$ iff $J_{W'}\subseteq I'$.
As every closed ideal is the intersection of the open ideals containing it, the
claim follows.

In particular $J_{W'}=0$ iff $\phi(J_W)R'=0$ iff $\phi$ factors through
$R/J_W$. Applying this to the universal deformation $W^{\univ}$ over
$R^{\univ}_{\rhobar}$ and $J:=J_{W^{\univ}}$: a morphism $\phi\colon
R^{\univ}_{\rhobar}\to R'$ has its associated deformation satisfying $P$ iff
$\phi$ factors through $R^{\univ}_{\rhobar}/J$. Hence $P_{\rhobar}$ is
represented by $R^{\univ}_{\rhobar}/J$.
\end{proof}

\begin{proposition}\label{prop:hr-condition}
For $k$ finite of characteristic $\ell\geq3$ and $\rhobar$ an absolutely
irreducible hardly ramified $k$-representation of $\Gal(\Q^{\{2,\ell\}}/\Q)$,
the condition ``hardly ramified'' satisfies \textup{(1)}--\textup{(3)}. Hence
there is a \emph{universal hardly ramified deformation ring}
$R^{\hr}_{\rhobar}$ and a universal hardly ramified deformation
$\rho^{\hr}\colon\GQ\to\GL_2(R^{\hr}_{\rhobar})$.
\end{proposition}

\begin{proof}[Proof sketch]
``Unramified outside $2\ell$'' is built into $G=\Gal(\Q^{\{2,\ell\}}/\Q)$.
``Cyclotomic determinant'' and ``flat at $\ell$'' satisfy (1)--(3) directly: (1)
is the definition of flatness, and (2),(3) follow because flatness is preserved
under finite sums and subquotients. The only condition needing care is (HR4) at
$2$, namely $\rho|_{\GQtwo}\cong\left(\begin{smallmatrix}\chi\alpha&*\\0&\alpha\end{smallmatrix}\right)$
with $\alpha$ unramified, $\alpha^2=1$. The forward implications of (1),(3) (that
the shape is preserved by quotients and base change) are clear. For the reverse
implications and (2), the key point is that the stable rank-one submodule is
\emph{unique}: since $\rho(\Frob_2)=\left(\begin{smallmatrix}2u&?\\0&u\end{smallmatrix}\right)$
with $u=\pm1$, its two eigenvalues $2u,u$ are distinct (their difference is
$u\in R^\times$), so it has a unique eigenline. This forces the following
descent lemma, which yields (2) and the reverse of (1),(3).

\emph{Lemma.} Let $R$ be a coefficient ring, $G$ a group acting on a free
rank-two $R$-module $W$, and $\alpha,\beta\colon G\to R^\times$ characters with
$\alpha\not\equiv\beta\pmod{\m_R}$. Let $R\subseteq S$ and suppose $W\otimes_RS$
has a free rank-one $G$-stable submodule with character $\beta$ and quotient
character $\alpha$. Then $W$ itself has such a submodule.

\emph{Proof.} Choose $g\in G$ with $\alpha(g)\not\equiv\beta(g)\pmod{\m_R}$, so
$u:=\beta(g)-\alpha(g)\in R^\times$. By Cayley--Hamilton
$\rho(g)$ has characteristic polynomial $(X-\alpha(g))(X-\beta(g))$, so
$e:=(\rho(g)-\alpha(g))/u$ satisfies $e(e-1)=0$; thus $e$ is an idempotent and
$W=eW\oplus(1-e)W$. Reducing mod $\m_R$, $eW/\m$ and $(1-e)W/\m$ are the
$\beta(g)$- and $\alpha(g)$-eigenspaces of $\rhobar(g)$, each of dimension one
(distinct eigenvalues), so by Nakayama $eW$ and $(1-e)W$ are free of rank one.
Finally, over $S$ we have $W\otimes S=(eW\otimes S)\oplus((1-e)W\otimes S)$ with
$g$ acting by $\beta(g),\alpha(g)$, so $eW\otimes S$ is the given $\beta$-stable
submodule; hence $eW=(eW\otimes S)\cap W$ is $G$-stable, as required. $\square$

Applying the descent lemma with the unramified character $\psi$ as quotient
character and $\chi\psi$ as submodule character --- distinct modulo $\m$ because
their ratio is $\chi$, with $\chibar(\Frob_2)=2\neq1$ in $k$ (well defined as
$\chi$ is unramified at $2$, since $\ell\neq2$) --- gives the reverse
implications and (2). Thus ``hardly ramified'' is a deformation condition, and
Theorem~\ref{thm:condition} supplies $R^{\hr}_{\rhobar}$.
\end{proof}

We shall also use, following Mazur's methods (and global class field theory),
the following, which we do not prove.

\begin{proposition}[assumed]\label{prop:krull}
If $\rhobar$ is an absolutely irreducible hardly ramified $k$-representation
with $\ell\geq5$, then $R^{\hr}_{\rhobar}$ is Noetherian and has Krull dimension
at least $1$.
\end{proposition}

\section{The automorphy lifting theorem}

To connect with the automorphic side we deform over a totally real field, and
must allow more ramification than ``hardly ramified''.

\begin{definition}\label{def:Smin}
Let $F$ be totally real, $\ell$ a prime unramified in $F$, $S$ a finite set of
finite places of $F$ not dividing $\ell$, and $\rhobar\colon\GF\to\GL_2(k)$
absolutely irreducible ($k$ finite of characteristic $\ell$). Call $\rhobar$
(and its deformations) \emph{$S$-minimal} if: $\rhobar$ is unramified outside
$\ell S$; $\det\rhobar=\chi$; $\rhobar$ is flat at every prime of $F$ above
$\ell$; and $\tr\rhobar(i)=2$ for every $i$ in an inertia group $I_P$ with
$P\in S$.
\end{definition}

By the same methods as Proposition~\ref{prop:hr-condition}, ``$S$-minimal'' is a
deformation condition, represented by a ring $R^{S\text{-}\minl}_{\rhobar}$.

\begin{example}
For $F=\Q$, $S=\{2\}$ and $\ell\geq5$, a hardly ramified representation is
exactly a $\{2\}$-minimal one: the condition $\tr\rhobar(i)=2$ on $I_2$ is
equivalent to (HR4), the hypothesis $\ell\geq5$ being needed to exclude
tamely ramified lower-triangular deformations at $2$. So $R^{\hr}_{\rhobar}$ is
a special case of a universal $S$-minimal deformation ring.
\end{example}

We can now state the ``Final Boss'', the one deep theorem of the twentieth
century that the whole endgame rests on. It is an ``$R=T$'' theorem: the
deformation ring equals a Hecke algebra. We do not prove it here (that is the
content of a sequel); we assume it and call it \textup{ALT}.

\begin{theorem}[Automorphy lifting theorem; Wiles, Kisin, \dots]\label{thm:alt}
Let $\ell\geq5$, let $F$ be totally real and unramified at $\ell$, let $k$ be
finite of characteristic $\ell$, and let $\rhobar\colon\GF\to\GL_2(k)$ be
absolutely irreducible, $S$-minimal. Assume in addition:
\begin{itemize}
\item $\rhobar|_{\Gal(\Fbar/F(\zeta_\ell))}$ is absolutely irreducible; and
\item $\rhobar$ is automorphic, arising from an eigenform in
$\calA_D(U_0(S);k)$ for a totally definite quaternion algebra $D/F$ split at all
finite places.
\end{itemize}
Then
\begin{enumerate}
\item[\textup{(a)}] $R^{S\text{-}\minl}_{\rhobar}$ is a module-finite
$\Zl$-algebra; and
\item[\textup{(b)}] if $K$ is a finite extension of $\Ql$ and
$R^{S\text{-}\minl}_{\rhobar}\to K$ is a homomorphism, the induced
representation $\rho\colon\GF\to\GL_2(K)$ arises from a system of $K$-eigenvalues
for $(F,D,S)$; that is, $\rho$ is automorphic.
\end{enumerate}
\end{theorem}

\section{The boss: $B_9$ implies FLT}

\begin{theorem}[Boss of Level 9]\label{thm:boss9}
Statements $B_{8a}$, $B_{8b}$, $B_{8c}$ together with the automorphy lifting
theorem imply Fermat's Last Theorem, modulo results known in the 1980s.
\end{theorem}

\begin{proof}
The hypotheses include $B_{8a}\wedge B_{8b}\wedge B_{8c}=B_8$, which by the boss
of the previous level implies FLT. Adjoining the automorphy lifting theorem only
strengthens the hypotheses, so $B_9\Rightarrow{}$FLT.
\end{proof}

\begin{remark}
The boss is trivial by design. The point of introducing the automorphy lifting
theorem now is that it is powerful enough to eat the other three heads
single-handedly. In Level 10 we shall see that it implies $B_{8a}$: the
universal hardly ramified ring $R^{\hr}_{\rhobar}$, being a quotient of an
$S$-minimal ring, becomes module-finite over $\Zl$, and, having Krull dimension
at least one (Proposition~\ref{prop:krull}), therefore admits a homomorphism to
the integers of a finite extension of $\Ql$ --- and out of that homomorphism
falls the desired characteristic-zero lift. Levels 11 and 12 then dispatch
$B_{8b}$ and $B_{8c}$.
\end{remark}

\end{document}
